\section{The Boundaries the Compiler Cannot See}
\label{sec:layer-boundaries}

The previous section said what belongs in each crate. This section is about what happens when the
answer is written down, agreed, and then quietly violated — because \emph{the compiler has no idea
that a layer boundary exists}. Every example in this section compiled. Every test passed. Every
lint was clean. That is what makes this class of defect expensive: there is no error message, only a
slow loss of the property the layering was supposed to provide.

\subsection{The Question That Finds Them}

There is a single question that finds these defects, and it is cheap to ask:

\begin{quote}
    \textbf{For this public signature, what does it force the caller to own?}
\end{quote}

A function that takes a register block forces the caller to have a handle to that register block.
A function that takes the whole peripheral singleton forces the caller to hold the singleton. Both
compile. Both are wrong the moment the caller is a layer above: the caller now has to reach for
something the layer below was supposed to keep to itself.

A real example, from this project. The HAL's buzzer and PWM constructors took the peripheral
singleton:

\begin{minted}{rust}
// Before: the caller has to own the PAC to start a tone.
pub fn new(dp: &pac::Peripherals, port: Port, pin: u8) -> Option<Self>
pub fn take(dp: &pac::Peripherals, port: Port, pin: u8) -> Option<Self>
// ... (bodies omitted: the listing is about the signatures)
\end{minted}

The signature looks innocent — it is the idiom of every Rust HAL, and it is how exclusivity is
proven at compile time. The consequence only showed up two layers up, in an example that had
already called \codeinline{Board::take()}:

\begin{minted}{rust}
// The example could not compile without this line, and nobody wrote down why.
let pac = unsafe { efr32mg24_pac::Peripherals::steal() };
let timer = Timer0::take(&pac, Port::A, 7);
\end{minted}

An example that has to steal the singleton is not an example with a problem; it is a HAL that
failed to hand out what it already owned. The fix is not to teach the example to steal more
politely. The fix is a door on the layer that owns the thing:

\begin{minted}{rust}
// After: the board owns the PAC once, and hands out what a caller needs.
let mut board = Board::take().expect("peripherals taken twice");
let mut led = board.pwm_on(Port::A, 7, 1_000)?;   // no PAC in sight
let mut mic = board.converter().microphone()?;    // a handle, not a register block
\end{minted}

The two constructs that make this work are worth naming, because they generalise:

\begin{itemize}
    \item \textbf{A constructor on the owner.} \codeinline{Board::pwm_on} is one line: it calls the
          private \codeinline{Timer0::take} with the singleton the board already holds. The
          information did not have to travel; it was already below.
    \item \textbf{An opaque handle.} \codeinline{microphone::Converter} is a \codeinline{Copy}
          wrapper around the register block, with a private constructor and one public method,
          \codeinline{microphone()}. The driver takes the handle; the caller never sees the type
          inside it.
\end{itemize}

\subsection{The Same Defect, One Layer Down}

Once you have the question, it stops being about the PAC. Three examples in the same project
reached for a \emph{core} peripheral — SysTick for a delay, DWT for the cycle counter:

\begin{minted}{rust}
let mut delay = Delay::with_sysclk(unsafe { cortex_m::Peripherals::steal().SYST }, 19_000_000);
let mut core = unsafe { cortex_m::Peripherals::steal() };
core.DCB.enable_trace();
core.DWT.enable_cycle_counter();
\end{minted}

Exactly the same shape: the layer that owns SysTick is the board, and it was not offering a door.
Worse than a style problem — two owners of SysTick is a delay that lies, and a cycle counter that
was never enabled reads as zero, which looks exactly like a measurement of a very fast loop. This
project lost a full round of fall-detection data to that second one: the loop declared 400 Hz and
ran at 158 Hz, and the only clue was that four different drops all showed the same 0.62 s of calm.

The fix, again, is doors rather than discipline:

\begin{minted}{rust}
let mut delay = board.delay();          // SysTick, handed out by its owner
let start = unsafe { dwt::start() };    // the cycle counter, switched on in one place
let elapsed = dwt::since(start);
\end{minted}

\subsection{A Layer That Exports the Layer Below}

The most damaging version of the same mistake is a re-export, because it does not look like code
at all:

\begin{minted}{rust}
// in a board support package
pub use efr32mg24_pac as pac;
\end{minted}

One line, entirely reasonable-looking, and it means every application can reconfigure every
register the board claimed to own. The board's API becomes advisory: the pins it configured, the
power gates it enabled, the clock tree it chose — all of it can be undone by a caller that has the
PAC. A BSP that re-exports its PAC is not a BSP; it is a suggestion.

The same crate carried a second, subtler version: \emph{two facades at once}. A newer one
(\codeinline{Module}, the LED and the on-board sensors) and an older one (\codeinline{Board} with
\codeinline{Pins} and a type-state GPIO layer) both existed, both public, and the examples had
split between them. Half the code proved one, half proved the other, and nothing exercised the
older one — which is the failure mode that matters: \textbf{an API with zero examples passes every
check trivially}. It compiled, it was documented, and it had been dead for weeks.

\subsection{Why It Kept Happening: The Table}

Seven defects of this class were found in a single pass over a codebase that built cleanly. What
they have in common is the column that says why nothing complained.

\begin{table}[htbp]
    \centering
    \footnotesize
    \begin{tabular}{@{}p{0.30\textwidth}p{0.30\textwidth}p{0.32\textwidth}@{}}
        \toprule
        \textbf{Defect} & \textbf{Why it compiled} & \textbf{What it cost} \\
        \midrule
        HAL constructors took a singleton or a register block
            & The signature is legal Rust, and idiomatic
            & examples had to \codeinline{steal()} the PAC \\
        BSP re-exported the PAC
            & A \codeinline{pub use} is invisible in review
            & the board's promises became advisory \\
        Two facades in one BSP crate
            & Both compiled; neither was unused by the compiler's standards
            & half the examples proved a dead API \\
        \codeinline{Board::gpio()} was public
            & Returning a register block is legal
            & callers reconfigured pins the board owned \\
        Two GPIO implementations in one HAL
            & Different modules, both public, no overlap warning
            & 605 unproven lines, and docs naming the wrong pin \\
        Documentation outlived the code
            & Prose is not compiled
            & 78 MHz, 256 KB, and an API that did not exist \\
        Examples stole SysTick and DWT
            & \codeinline{steal()} is \codeinline{unsafe} but compiles
            & a delay that lied, a counter that read zero \\
        \bottomrule
    \end{tabular}
    \caption{Seven layer violations in a codebase where cargo build, cargo test and cargo clippy
    were all green.}
    \label{tab:layer-defects}
\end{table}

\subsection{Writing the Boundary Down as a Check}

A boundary that is only a paragraph is a boundary that will be crossed, and this project crossed
all seven. The rules only started holding when each one became a line in the gate script — the same
script that checks formatting and tests, so that it runs when nothing else does:

\begin{minted}{bash}
# 1. An example may only name the layer it demonstrates.
grep -rln 'efr32mg24_pac|pac::' crates/*/examples/ | grep -v '^crates/efr32mg24-pac/'
# 2. A BSP must not re-export the PAC.
grep -rln 'pub use efr32mg24_pac' crates/*/src/ | grep -v '^crates/efr32mg24-hal/'
# 3. The HAL must not put a register block in a public signature.
grep -rqE 'pub (fn|const fn) [^\n]*::RegisterBlock' crates/efr32mg24-hal/src/
# 4. No example may steal a core peripheral.
grep -rln 'Peripherals::steal' crates/*/examples/
# 5. Every crate carries at least one example of its own level.
for dir in crates/*/ libraries/*/; do find "$dir/examples" -maxdepth 1 -name '*.rs'; done
\end{minted}

Each of those five checks failed at least once during the pass that wrote them. That is the
argument for having them: not that the code was wrong, but that the code was wrong in ways that
had already survived review.

\subsection{What Is Not a Violation}

The distinction that took the longest to state clearly is between a \emph{token} and a
\emph{register block}:

\begin{minted}{rust}
// Legitimate: the singleton token proves exclusivity at compile time.
pub fn new(usart: pac::Usart0S, config: Config, clocks: &FrozenClocks) -> Self

// A violation: the caller receives the registers themselves.
pub fn new(adc: &pac::iadc0_s::RegisterBlock) -> Result<Self, MicError>
// ... (bodies omitted: the listing is about the signatures)
\end{minted}

Both mention \codeinline{pac}. The first is how every Rust HAL on the market works —
\codeinline{dp.USART0} in \codeinline{stm32f4xx-hal}, \codeinline{peripherals.USART0} in
\codeinline{rp2040-hal} — and it is the mechanism that makes it impossible to create two USART
drivers for one USART. The second hands over the hardware itself, with no ownership at all: any
number of callers may now configure the same converter, and the driver's invariants are whatever
the last writer decided.

The rule, then, is not ``never mention the PAC in the HAL''. It is: \textbf{the HAL may take a
token, and may never give out registers}.

\subsection{The Procedure, for the Next Project}

The pass that found seven defects was mechanical, and it transfers:

\begin{enumerate}
    \item \textbf{List every public signature} in every crate, oldest and least used first. The dead
          API is where the violations accumulate, because nothing compiles against it.
    \item \textbf{Ask the question} of each one — what does this force the caller to own? — and
          write the answer down next to the signature.
    \item \textbf{Close each leak with a door on the owner}: a method on the struct that already
          holds the thing, or an opaque handle for what must be passed through.
    \item \textbf{Turn each rule into a check} in the gate that already runs, and confirm the check
          fails on the old code before accepting it. A check that has never been red has not been
          tested.
\end{enumerate}

The last step is the one that decides whether the work survives. This project had already written
the layering down — in the README, in this book, in the crate documentation — and had violated it
seven times anyway. What changed the outcome was not a better description of the layers. It was
five greps that fail the build.

\begin{keytakeaways}
    \item A layer boundary is invisible to the compiler: the code compiles either way, which is why
          this class of defect survives review.
    \item The question that finds them is \emph{what does this force the caller to own?}
    \item A HAL may take a peripheral \emph{token} to prove exclusivity, and may never hand out
          \emph{registers} — the difference decides whether a driver's invariants mean anything.
    \item A BSP that re-exports the PAC makes its own configuration advisory.
    \item An API with zero examples passes every check trivially; it is dead and nobody notices.
    \item A rule that is only prose will be broken. A rule that is a line in the gate script is a
          rule that holds.
\end{keytakeaways}
